\begin{thebibliography}{100}
\bibitem{1}経済産業省・資源エネルギー庁：”第2部エネルギー動向第１章国内エネルギー動向第1説エネルギー需供の動向”，\\https://www.enecho.meti.go.jp/about/whitepaper/2023/pdf/2\_1.pdf，閲日時2025年2月20日17:20.
\bibitem{2}環境省：「令和3年度版環境白書循環型社会白書/生物多様性白書」
\bibitem{3}河村篤男，”現代パワーエレクトロニクス”，数理工学社，2005
\bibitem{4}角田秀夫，”PLLの基本と応用”，東京電機大学出版局，1978
\bibitem{5}石原清志，河村篤男："DSPによる実時間出力合成形三相PWMインバータ"電気学会D，110巻6号，平成2年
\bibitem{6}川島加也：”Universal Smart Power Module(USPM)コントローラを用いた系統連系システム”，令和3年度東京電機大学卒業論文
\bibitem{7}加藤遼子：”系統連系用マルチレベルインバータにおける外乱補償型デッドビート制御系の構築”，令和3年度東京電機大学修士論文
\bibitem{8}廣惠大輔：”高性能FPGAコントローラによる50MHzマルチサンプリングデッドビート制御”，令和2年度東京電機大学卒業論文
\bibitem{9}細山田悠：”高速High/LOWパルス動作を実現する3相インターリーブDC/DCコンバータの制御手法の検討”，令和3年度横浜国立大学博士論文
\bibitem{10}小原秀嶺：”フライングキャパシタマルチレベルコンバータによる電力変換器の波形改善および高パワー密度化に関する研究”，千葉大学2015年度博士論文
\bibitem{11}窪慎之助：”連携リアクトルなし三相系統連系システムのロバスト性検証”，令和4年度東京電機大学卒業論文
\bibitem{12}中村一稀：”SoC-FPGA によるUniversal Smart Power Module(USPM)コントローラを用いた三相系統連系インバータの12.5MHz マルチサンプリング外乱補償型デッドビート制御”，令和4年度東京電機大学修士論文
\bibitem{13}白崎圭亮, 平川遼太郎, 天野博之, 父母靖二, 片岡良彦: “ 系統安定化機能付き GFL インバータの使用設計 ”, 電気学会全国大会, p403, 2024.
\bibitem{14}高野明夫: ”TNPC マルチレベルインバターのVFM 制御”, 沼津高等専門学校研究報告. Vol.56, 2022.
\bibitem{15}浦壁隆浩, 西尾直樹, 藤原賢司, 川上知之: “ 太陽光発電システム-高効率パワーコンディショナの技術- ”, 三菱電機技報, 2009.
  %\bibitem{h}
  %\bibitem{i}
  %\bibitem{l}
  %\bibitem{m}
  %\bibitem{n}
  %\bibitem{o}
  %\bibitem{p}
  %\bibitem{q}
  %\bibitem{r}
  %\bibitem{s}
  %\bibitem{t}
  %\bibitem{u}
\end{thebibliography}